\documentclass[10pt]{article}
\usepackage[margin=0.56in]{geometry}
\usepackage{amsmath,amssymb,booktabs,graphicx,microtype,xcolor}
\usepackage[hidelinks]{hyperref}
\setlength{\parindent}{0pt}
\setlength{\parskip}{3pt}
\setlength{\abovedisplayskip}{4pt}
\setlength{\belowdisplayskip}{4pt}
\setlength{\textfloatsep}{5pt}
\renewcommand{\arraystretch}{1.08}

\newcommand{\WarmPhaseAgreement}{87.5\%}
\newcommand{\JointPhaseAgreement}{37.5\%}
\newcommand{\WarmGateAgreement}{91.2\%}
\newcommand{\JointGateAgreement}{49.1\%}
\newcommand{\WarmShareMAE}{0.152}
\newcommand{\JointShareMAE}{0.337}
\newcommand{\ProtocolDisagreement}{59.3\%}
\newcommand{\CompositionalRuns}{143}
\newcommand{\TotalRuns}{432}
\newcommand{\CappedRuns}{333}

\title{\vspace{-1.15cm}\textbf{Two-Block Circuit Formation:\\Stationary Theory versus Training}}
\author{}
\date{\vspace{-1.25cm}}

\begin{document}
\maketitle
\vspace{-0.7cm}

\textbf{Result in one sentence.}
The fully trainable network can form a $2+2\!\to\!4$ circuit, but only in a finite
layerwise-decay window: weak first-block decay leaves a cheap block-1 shortcut, whereas strong
first-block decay destroys the degree-2 parents; initialization determines whether training
finds the compositional minimum inside that window.

\textbf{Theory tested.}
For singleton norm $s=Q_1^{(0)}$, parent norm $q=Q_2$, and
$Q_1^{(1)}=\sqrt{s^2+q^2}$, the reduced zero-temperature theory compares the optimized
direct hybrid $H_D=D_1+D_2$ with the compositional hybrid $H_C=D_1+C$.  At fixed
representation, block 2 selects $C$ precisely when
\[
R\equiv Q_2/Q_1^{(1)}>1/\sqrt8.
\]
The representation is itself optimized, so the actual boundary follows from $F_C=F_D$.
Inside $H_C$, the predicted compositional output share is
\[
f_C=\frac{8\lambda_1q^2}{\lambda_2s^2+8\lambda_1q^2}.
\]
Increasing $\lambda_1/\lambda_2$ therefore suppresses the block-1 shortcut but also raises the
cost of the block-1 parents, producing a bounded rather than monotone compositional region.

\textbf{Experiment.}
We trained a bias-free two-residual-block ReLU network of width $N=1024$ on the exact
four-bit cube.  Both degree-2 parents and the degree-4 target were trained simultaneously,
with no frozen weights and
$\gamma_2=\gamma_4=2$.  The tested plane was
\[
\lambda_E=\lambda_2=\lambda_p=\lambda_4=\lambda_0,
\qquad \lambda_1=(\lambda_1/\lambda_2)\lambda_0,
\]
using nine $\lambda_0$ values, eight ratios, three seeds, and two protocols.  \emph{Joint}
trains both parent targets and the degree-4 target from the common random initialization,
with every weight trainable.  \emph{Parent-warm} first trains $E,V_1,W_{d2}$ on the parents
read from $h_1$, then unfreezes everything and trains the identical final objective
(216 runs per protocol; \TotalRuns{} total).  Full-batch float64 SGD used
$\eta_{E,V}=10^{-2}$ and the muP readout scaling $\eta_W=10^{-2}/N$.  Each cell ran for at
least 15k and at most 50k steps.  Every empirical row below uses the protocol specified
before training; no post-hoc selection between protocols is used.

\begin{figure}[!ht]
\centering
\includegraphics[width=0.66\textwidth]{../analysis_output/two_block_phase_diagram/theory_empirical_comparison.pdf}
\vspace{-7pt}
\caption{\textbf{Theory (top), parent-warm (middle), and joint training (bottom).}
Columns show phase, total compositional target share, and final accessibility.  Cyan marks
$R=1/\sqrt8$.  Each empirical row averages its own three seeds; neither row is selected or
replaced using the other protocol's final objective.}
\end{figure}

\clearpage
\textbf{Quantitative comparison.}

\begin{center}
\begin{tabular}{@{}lcc@{}}
\toprule
Diagnostic & Parent-warm & Joint \\
\midrule
Theory/protocol phase agreement over the $8\times9$ grid
  & \WarmPhaseAgreement & \JointPhaseAgreement \\
Measured-$R$ gate agreement in active cells
  & \WarmGateAgreement & \JointGateAgreement \\
Mean absolute error in total compositional share $f_C$
  & \WarmShareMAE & \JointShareMAE \\
Grid cells with majority phase $H_C$ & $44/72$ & $3/72$ \\
Individual runs classified $H_C$ & $134/216$ & $9/216$ \\
\bottomrule
\end{tabular}
\end{center}

The two protocols select different phases in \ProtocolDisagreement{} of matched
seed/cell comparisons (128/216).

\textbf{What the phase diagram establishes.}
First, the desired circuit is feasible without freezing: the prespecified parent-warm
protocol occupies a substantial band at moderate $\lambda_0$ and moderate
$\lambda_1/\lambda_2$, and reproduces the theoretical phase topology much better than joint
training.  Second, the two
failure modes predicted by the reduced theory are distinct.  At small ratio, parents remain
large but the degree-4 readout uses the inexpensive $D_1$ shortcut.  At large ratio, $D_1$
is suppressed but $Q_2$ collapses, $R$ falls below $1/\sqrt8$, and block 2 uses $D_2$ or the
target is suppressed.  Third, matched joint and parent-warm runs---same seed, same decay cell,
and identical final objective---selected different phases in 128 of 216 cases.  This
hysteresis is finite-time/nonconvex basin selection, not a shift of the local gate boundary.
The largest parent-warm mean compositional share was $0.681$ at $\lambda_0=10^{-3}$ and
$\lambda_1/\lambda_2=1$, with $R=0.595$ and target amplitude $0.995$; the theory selects
$H_C$ in the same cell and predicts $f_C=0.749$.

\textbf{How to favor $2+2\!\to\!4$.}
Use nonzero decay on the degree-4 readout so hidden/readout rescaling cannot evade the norm
comparison.  Keep $\lambda_E$, $\lambda_p$, and especially the parent-producing part of
$\lambda_1$ small enough that the \emph{final} measured $R$ stays above $1/\sqrt8$; then make
block 2 cheaper than the block-1 degree-4 shortcut by lowering $\lambda_2$ relative to the
target-specific block-1 decay.  Whole-layer $\lambda_1$ is an imperfect control because it
taxes wanted parents and the unwanted shortcut together.  If the feasible window is narrow,
separate those directions with a decayed block-1 degree-4 adapter or soft subspace penalty.

\textbf{Interpretation and limits.}
The stationary theory captures the topology of the control problem---a compositional window
between shortcut domination and parent collapse---but quantitative differences are expected.
The finite network contains mixed Walsh features, non-orthogonal singleton leakage, and
multifunctional neurons omitted by the balanced rank-one ansatz.  Moreover, many runs hit
the 50k-step cap (\CappedRuns{} of \TotalRuns{}), while protocol disagreement directly
demonstrates metastability.  Thus the parent-warm map is a history-conditioned accessibility
map, not proof of equilibrium, while the joint map directly displays the optimization
barrier from random initialization.  The most robust local diagnostic remains the measured inequality
$Q_2/Q_1^{(1)}>1/\sqrt8$, together with the signed target contributions from $h_1$ and the
second residual branch.

\begin{figure}[h]
\centering
\includegraphics[width=0.52\textwidth]{../analysis_output/two_block_phase_diagram/full_objective_protocol_difference.pdf}
\vspace{-7pt}
\caption{\textbf{Full regularized objective difference between the two prespecified
protocols.}  Each cell shows the mean matched-seed difference
$J_{\rm warm}-J_{\rm joint}$ using a signed-log color scale: blue favors parent-warm, red
favors joint, and the black contour marks equality.  This is an objective diagnostic only;
it is not used to select or replace either empirical phase-map row in Fig.~1.}
\end{figure}

\end{document}
